\section{Проверка реализации и воспроизводимость}
\subsection{Автоматические проверки}
В актуальном окружении выполнена сборка и запуск всех четырёх целей CTest. Результат: 4 из 4 наборов прошли. Это относится к совокупности наборов, а не к числу отдельных тестовых случаев.
\begin{longtable}{L{0.3\linewidth}L{0.62\linewidth}}
\toprule Цель & Проверяемые свойства \\\midrule\endhead
\code{calendarservice} & Входные данные, права, конфликты ресурсов, интервалы и границы, фильтры, жизненный цикл и архив. \\
\code{simulationengine} & Ежедневные планёрки всех отделов, семинары, состав руководителей, короткий срок планирования, генерация и независимость от шага. \\
\code{completion} & Напоминания, дедупликация и версии, сохранение/продолжение, метрики, экспериментальные агрегаты, явный отказ для прежних генераторов. \\
\code{calendarwindow} & Передача параметров, календарные формы, права ролей, ручной выбор переноса, возврат к заявке, фильтры, пауза и визуальное движение. \\
\bottomrule
\end{longtable}
Логи и JUnit-результат сохранены в каталоге:\begin{quote}\file{documentation/experiments/verification-v7}\end{quote}Дополнительно скрипт генерации таблиц повторно вычисляет показатели исследования из отдельных прогонов. Успешные автоматические проверки подтверждают проверенные случаи, но не заменяют проверку всех возможных пользовательских сценариев и согласование требований преподавателем.

\subsection{Воспроизведение исследования и сборка отчёта}
Для повторения серии нужны та же версия алгоритма и совместимая реализация стандартных случайных распределений C++. Разные стандартные библиотеки могут дать другую случайную последовательность даже при одном seed; идентичность между платформами не обещается.
\begin{lstlisting}
cmake -S . -B build/qt6 -DBUILD_TESTING=ON
cmake --build build/qt6 -j2
ctest --test-dir build/qt6 --output-on-failure
./build/qt6/modeling_experiments \
  --output documentation/experiments/results-v7 \
  --repetitions 20
python3 documentation/report/make_results.py
cd documentation/report
xelatex -interaction=nonstopmode -halt-on-error report.tex
xelatex -interaction=nonstopmode -halt-on-error report.tex
\end{lstlisting}
Альтернатива двум вызовам XeLaTeX --- \code{tectonic report.tex}. Для диаграмм используется TikZ, для кириллицы --- fontspec и babel. Основной шрифт --- DejaVu Serif; если он отсутствует, используется Times New Roman. Аналогично предусмотрены запасные системные шрифты для гротеска и моноширинного текста.

\subsection{Соответствие пунктам отчёта из задания}
\begin{longtable}{L{0.65\linewidth}L{0.27\linewidth}}
\toprule Требование & Раздел \\\midrule\endhead
Уточнённая постановка задачи & 1 \\
Диаграмма классов & 2 \\
Диаграмма объектов в рамках проектирования & 2 \\
Текстовые спецификации основных классов & 3 \\
Инструментальные средства & 4 \\
Файловая структура & 5 \\
Краткая характеристика интерфейса & 6 \\
Описание проведённых экспериментов & 7 \\
Подготовка полного Git-репозитория на GitHub & 8: указан адрес и фактический статус публикации \\
\bottomrule
\end{longtable}
Открытые вопросы: согласование имитации электронной почты и публикация актуального комплекта на GitHub. Эти действия не подменены описанием и не заявлены выполненными.
